\documentclass[11pt,a4paper]{article}
\usepackage[utf8]{inputenc}
\usepackage[T1]{fontenc}
\usepackage{lmodern}
\usepackage{geometry}
\usepackage{hyperref}
\usepackage{graphicx}
\usepackage{booktabs}
\usepackage{url}
\geometry{margin=1in}
\hypersetup{colorlinks=true, linkcolor=blue, urlcolor=blue, citecolor=blue}
\providecommand{\RawCorpusLines}{71876}
\providecommand{\DedupDocs}{42217}
\providecommand{\SummaryDocs}{42217.0}
\providecommand{\SummaryVocab}{214003.0}
\providecommand{\SummaryTokens}{37491119.0}
\providecommand{\SummaryAvgLen}{888.0573939408296}
\providecommand{\BMNDCGTen}{0.8300558702936255}

\title{Temple University Information Retrieval Project\\
A Reproducible Search System for a Domain-Specific Web Collection}
\author{Project Report}
\date{\today}

\begin{document}
\sloppy
\maketitle
\tableofcontents
\bigskip

\section{Abstract}
This report presents an end-to-end information retrieval pipeline built over a Temple University web crawl. The system covers the full retrieval workflow: web collection processing, document deduplication, indexing, collection analysis, query construction, relevance assessment, candidate pooling, and evaluation of standard retrieval baselines. We use the resulting collection to study the behavior of classical retrieval models on a realistic, domain-specific corpus. The report summarizes the final system artifacts, implementation choices, experimental setup, and key empirical findings.

\section{Introduction}
The goal of this project is to develop a reproducible information retrieval system and evaluate how established retrieval models perform on a real web-derived collection. Unlike controlled benchmark datasets, web crawls contain duplicate pages, noisy HTML content, heterogeneous document lengths, and inconsistent metadata. These properties make the collection useful for studying practical retrieval challenges in a realistic setting.

Our corpus consists primarily of HTML pages crawled from the Temple University domain. After preprocessing and deduplication, the documents are indexed using PyTerrier/Terrier and evaluated with a set of annotated information needs. Where available, graded relevance judgments are used to compare retrieval effectiveness across multiple ranking models.

\section{Data Collection and Normalization}
The raw collection was stored as JSONL records and then normalized into a deduplicated corpus suitable for indexing and evaluation. The main collection statistics are:
\begin{itemize}
  \item Raw crawl records: \textbf{\RawCorpusLines} (\texttt{data/raw/corpus.jsonl}).
  \item Deduplicated documents: \textbf{\DedupDocs} (\texttt{data/raw/corpus\_deduped.jsonl}).
\end{itemize}

Deduplication was performed by retaining the first occurrence of each \texttt{docno}. This step ensures that each indexed document has a stable identifier and prevents duplicate records from artificially influencing retrieval results. The normalized corpus therefore provides a consistent basis for indexing, ranking, and evaluation.

\section{Indexing}
The deduplicated corpus was indexed using PyTerrier/Terrier in \texttt{index/terrier}. Each document was represented with \texttt{docno} metadata and a tokenized text field extracted from the normalized page content. For this snapshot of the collection, the resulting index contains \textbf{\SummaryDocs} documents, a vocabulary of \textbf{\SummaryVocab} terms, \textbf{\SummaryTokens} total tokens, and an average document length of \textbf{\SummaryAvgLen} tokens.

These statistics provide a compact summary of the indexed collection and help contextualize the retrieval experiments. In particular, vocabulary size, token count, and average document length are important indicators of collection scale and lexical variability.

\section{Experimental Setup}
We evaluated several standard bag-of-words retrieval baselines, including BM25, TF--IDF, and Query Likelihood. All models were evaluated using graded relevance judgments when available. In addition to comparing baseline effectiveness, we conducted a BM25 parameter sensitivity analysis over $(k_1, b)$ to examine how ranking performance changes under different term-frequency saturation and length-normalization settings.

We also ran a small optimization experiment comparing conjunctive and disjunctive query evaluation strategies. This analysis helps characterize the trade-off between retrieval strictness and candidate coverage, especially for multi-term information needs.

\section{Challenges and Implementation Notes}
% How the pipeline works and what broke along the way
\section*{System overview and challenges}
\subsection*{End-to-end flow}
\begin{enumerate}
  \item \textbf{Crawl} (\texttt{src/crawl\_temple.py}): BFS on \texttt{temple.edu}, stream-append JSONL, resume by scanning existing \texttt{docno} and URLs. Two queues prefer unseen URLs before re-fetching known pages (expansion-only).
  \item \textbf{Dedupe} (\texttt{src/dedupe\_corpus.py}): first row per \texttt{docno} $\rightarrow$ \texttt{data/raw/corpus\_deduped.jsonl}.
  \item \textbf{Index} (\texttt{src/build\_index.py}): PyTerrier \texttt{IterDictIndexer}, meta \texttt{docno}, text field \texttt{text}.
  \item \textbf{Analysis} (\texttt{src/collection\_analysis.py}): vocabulary, Zipf, lengths, stopwords, word cloud $\rightarrow$ \texttt{outputs/analysis/}.
  \item \textbf{Queries / pooling}: templates and \texttt{queries.tsv}; pool from BM25, TF-IDF, QL.
  \item \textbf{Bootstrap qrels}: rank-1 BM25 labeled relevant (temporary metrics only).
  \item \textbf{Eval, sensitivity, optimization}: PyTerrier experiments $\rightarrow$ \texttt{outputs/eval/}, \texttt{outputs/optimization/}.
  \item \textbf{Reports}: \texttt{src/update\_reports.py} refreshes Markdown; \texttt{src/export\_latex\_reports.py} emits this PDF bundle.
\end{enumerate}

\subsection*{Snapshot numbers (from this export)}
\begin{itemize}
  \item Raw JSONL lines: \textbf{{\RawCorpusLines}} (includes duplicate lines from restarts).
  \item Deduped documents: \textbf{{\DedupDocs}}.
  \item Collection summary (first CSV row): docs {\SummaryDocs}, vocab {\SummaryVocab}, tokens {\SummaryTokens}, avg length {\SummaryAvgLen}.
  \item BM25 nDCG@10 (aggregate file, bootstrap qrels): \textbf{{\BMNDCGTen}}.
\end{itemize}

\subsection*{Main challenges encountered}
\begin{itemize}
  \item \textbf{Resume and duplicates:} repeated crawl runs appended lines; some URLs mapped to multiple \texttt{docno} values over time. Dedupe-by-\texttt{docno} plus index-time \texttt{drop\_duplicates} kept Terrier stable.
  \item \textbf{Frontier / scheduling:} treating ``already saved'' URLs as non-expandable starved BFS; fixed by always enqueueing unvisited links and then splitting \emph{new} vs \emph{expand} queues so new URLs are tried first.
  \item \textbf{Scale vs site reality:} many hosts return HTTP 403 to datacenter or script clients; tqdm only moves on \emph{new} docs so long runs looked idle while HTTP work continued---postfix counters (\texttt{fetched}, queue sizes) were added for visibility.
  \item \textbf{PyTerrier/Java:} JVM init, absolute index paths, \texttt{text\_attrs}, unique \texttt{docno} constraints required several integration passes.
  \item \textbf{Evaluation wiring:} PyTerrier \texttt{Experiment} metric tables differ in wide vs aggregate layouts; bootstrap qrels are not a substitute for human \texttt{qrels.tsv}.
\end{itemize}


\section{Results}
\subsection{Plots}

\begin{figure}[htbp]
\centering
\includegraphics[width=\linewidth,keepaspectratio]{\detokenize{../../outputs/analysis/zipf_loglog.png}}
\caption{Zipf distribution (log--log) with dashed OLS fit.}
\end{figure}

\begin{figure}[htbp]
\centering
\includegraphics[width=\linewidth,keepaspectratio]{\detokenize{../../outputs/analysis/doc_length_distribution.png}}
\caption{Histogram of tokenized document lengths.}
\end{figure}

\begin{figure}[htbp]
\centering
\includegraphics[width=\linewidth,keepaspectratio]{\detokenize{../../outputs/analysis/wordcloud_top_terms.png}}
\caption{Word cloud of frequent non-stopwords.}
\end{figure}

\begin{figure}[htbp]
\centering
\includegraphics[width=\linewidth,keepaspectratio]{\detokenize{../../outputs/eval/bm25_sensitivity_curves.png}}
\caption{BM25 parameter sensitivity (nDCG@10 contours / curves).}
\end{figure}

\begin{figure}[htbp]
\centering
\includegraphics[width=\linewidth,keepaspectratio]{\detokenize{../../outputs/charts/aggregate_model_comparison.png}}
\caption{Aggregate retrieval metrics comparison.}
\end{figure}

\begin{figure}[htbp]
\centering
\includegraphics[width=\linewidth,keepaspectratio]{\detokenize{../../outputs/charts/ndcg_by_query_type.png}}
\caption{nDCG@10 broken down by query type.}
\end{figure}

\begin{figure}[htbp]
\centering
\includegraphics[width=\linewidth,keepaspectratio]{\detokenize{../../outputs/charts/per_query_ndcg_distribution.png}}
\caption{Dispersion of per-query nDCG@10.}
\end{figure}


\subsection{Stopwords, queries, and optimization}
\subsection*{Summary: stopwords}

Token coverage uses NLTK English stopwords extended with corpus low-IDF terms (\texttt{corpus\_tfidf\_stopwords.csv} lists 250 corpus-specific candidates; the merged list has 448 types after union with NLTK). In a sample pass, 26.00\% of tokens matched NLTK-only; 27.12\% matched the TF--IDF low-IDF list only; after merging, 53.12\% of tokens were treated as stopwords (vocabulary size from the summary file: 214003.0 types).

\subsection*{Summary: queries}

The evaluation set contains 105 annotated queries in \texttt{data/processed/query\_templates.tsv} (exported to \texttt{queries.tsv} for retrieval). Intent types: informational (39), navigational (34), broad\_ambiguous (20), broad (10), transactional (2). Pooling and \texttt{qrels.tsv} judgments are described in the weekly sections; full tables remain under \texttt{outputs/eval/} if needed.

\subsection*{Summary: optimization}

The optimization step compared \textbf{disjunctive} vs \textbf{conjunctive} BM25 query evaluation. Mean batch latency was about 1929~ms (disjunctive) vs 1618~ms (conjunctive); documents scored were 105000 vs 92390. Overlap at rank 10 between the two rankings was approximately 0.520.

\subsection*{Metrics and models}

The figures above summarize sensitivity, query-type effects, and dispersion of per-query scores. The table below gives aggregate MAP, P@5, recall@10, and nDCG@10 for each run model. Use the exported per-query CSV for head-to-head contrasts beyond these totals.

\begin{table}[htbp]
\centering
\footnotesize
\caption{Aggregate retrieval metrics (graded qrels as in the eval run). PyTerrier baselines; optional Elasticsearch BM25 row when the extension was executed.}
\resizebox{\textwidth}{!}{%
\begin{tabular}{l r r r r}
\toprule
name & map & P\_5 & recall\_10 & ndcg\_cut\_10 \\
\midrule
BM25 & 0.4897853550849345 & 0.9961904761904762 & 0.4184918391876504 & 0.8300558702936255 \\
TFIDF & 0.5078014044068315 & 0.9942857142857143 & 0.4127164492343494 & 0.8243925857440955 \\
QL & 0.5411758282026168 & 0.937142857142857 & 0.3935576203971061 & 0.7732880292938707 \\
ElasticsearchBM25 & 0.6048140709232691 & 0.6895238095238095 & 0.25539404090465273 & 0.5027055516290783 \\
\bottomrule
\end{tabular}
}
\vspace{4pt}
\textit{\footnotesize{Sensitivity grids, per-query scores, and query-type splits are on disk under \texttt{outputs/eval/}; charts under \texttt{outputs/charts/}.}}
\end{table}


\subsection*{Formal hypothesis tests}

PyTerrier paired comparisons: \texttt{\detokenize{outputs/eval/paired_t_tests.csv}}. Elasticsearch vs BM25: \texttt{\detokenize{outputs/extension/es_vs_pyterrier_ttest.csv}}.

\section{Reproducibility}
All reported outputs can be regenerated from the project artifacts. To refresh the intermediate outputs and rebuild this PDF, run:
\begin{verbatim}
.venv/bin/python3 -m src.refresh_phase_outputs
cd reports/tex && pdflatex main.tex && pdflatex main.tex
\end{verbatim}

This workflow regenerates the phase outputs, updates the LaTeX macros, and compiles the final report from the source files.
\end{document}
